\documentclass[10pt]{article}
\usepackage[T1]{fontenc}
\usepackage[utf8]{inputenc}
\usepackage[french]{babel}
\usepackage{amsmath,amssymb,mathrsfs,mathpazo}
\usepackage{geometry}
\geometry{a4paper,margin=2cm}
\usepackage{graphicx,array,tabularx,tabularray,colortbl,multirow,multicol,booktabs,enumitem,pifont,ulem,fancybox,diagbox,cancel}
\usepackage{pstricks,pstricks-add,pst-plot,pst-tree,pst-node,pst-text,pst-func,pst-math,pst-bspline,pst-3dplot,pst-eucl,pst-solides3d}
\let\pscalculate\undefined
\usepackage{pst-bezier}
\usepackage{tikz,pgfplots,tkz-tab}
\usetikzlibrary{arrows,arrows.meta,patterns,calc,positioning,trees,decorations.markings,shapes.geometric,shapes.misc,shapes.arrows,intersections,decorations.pathmorphing,decorations.pathreplacing,backgrounds,shadows,babel}
\pgfplotsset{compat=1.18}
\usepgfplotslibrary{fillbetween}
\usepackage[np]{numprint}
\usepackage{listings,makecell,pas-tableur,twoopt}
\let\barre\undefined
\usepackage[pstricks]{tablvar}
\definecolor{darkgreen}{RGB}{0,100,0}
\providecommand{\red}{\color{red}}
\providecommand{\blue}{\color{blue}}
\providecommand{\green}{\color{green}}
\providecommand{\black}{\color{black}}
\providecommand{\white}{\color{white}}
\providecommand{\gray}{\color{gray}}
\providecommand{\cyan}{\color{cyan}}
\providecommand{\magenta}{\color{magenta}}
\providecommand{\yellow}{\color{yellow}}
\providecommand{\violet}{\color{violet}}
\providecommand{\orange}{\color{orange}}
\providecommand{\euro}{\texteuro}
\providecommand{\EUR}{\texteuro}
\providecommand{\ieme}{\textsuperscript{e}}
\providecommand{\degres}{\ensuremath{^\circ}}
\providecommand{\degree}{\ensuremath{^\circ}}
\providecommand{\vv}[1]{\overrightarrow{#1}}
\providecommand{\Coord}[1]{#1}
\providecommand{\decosix}{}
\providecommand{\starredbullet}{\ensuremath{\bullet}}
\pagestyle{empty}
\setlength{\parindent}{0pt}
\usepackage{adjustbox}\DeclareUnicodeCharacter{2236}{:}\DeclareUnicodeCharacter{2190}{\ensuremath{\leftarrow}}
\providecommand{\R}{\mathbb{R}}
\providecommand{\N}{\mathbb{N}}
\providecommand{\Z}{\mathbb{Z}}
\providecommand{\Q}{\mathbb{Q}}
\providecommand{\D}{\mathbb{D}}
\providecommand{\C}{\mathbb{C}}
\providecommand{\e}{{\mathrm{e}}}
\providecommand{\vect}[1]{\overrightarrow{#1}}
\providecommand{\vectt}[1]{\overrightarrow{\mathrm{#1}}}
\providecommand{\vvt}[1]{\overrightarrow{\mathrm{#1}}}
\providecommand{\np}[2][]{{#2~#1}}
\providecommand{\vv}[1]{{\overrightarrow{#1}}}
\renewcommand{\d}{{\mathrm{d}}}
\providecommand{\degres}{{^\circ}}
\providecommand{\degre}{{^\circ}}
\providecommand{\blue}{\color{blue}}
\providecommand{\red}{\color{red}}
\providecommand{\enskip}{\quad}
\providecommand{\llbracket}{\lbrack\!\lbrack}
\providecommand{\rrbracket}{\rbrack\!\rbrack}
\providecommand{\euro}{{\text{€}}}
\providecommand{\no}{{\text{n°}}}
\providecommand{\interval}[2]{{[#1;#2]}}
\providecommand{\intervFO}[2]{{[#1;#2[}}
\providecommand{\dotfill}{\ldots}

\providecommand{\ds}{\displaystyle}
\providecommand{\textding}[1]{#1}
\providecommand{\Oij}{\(\mathrm{O};\vec{\imath},\vec{\jmath}\)}
\providecommand{\Oijk}{\(\mathrm{O};\vec{\imath},\vec{\jmath},\vec{k}\)}
\providecommand{\Ouv}{\(\mathrm{O};\vec{u},\vec{v}\)}
\providecommand{\intervFF}[2]{\ensuremath{[#1;#2]}}
\renewcommand{\euro}{\ifmmode\text{\texteuro}\else\texteuro\fi}
\DeclareUnicodeCharacter{00B0}{\ifmmode^{\circ}\else\textdegree\fi}
\DeclareUnicodeCharacter{20AC}{\ifmmode\text{\texteuro}\else\texteuro\fi}
% Énoncés : Éducation nationale. Transcription et corrigés : APMEP, Denis Vergès et contributeurs. Découpage et classement : CoopMaths, Sébastien Lozano, Éric Elter et contributeurs. Adaptation de la présentation : Maths Explorer.
% CC BY-SA : https://coopmaths.fr/www/about/
% https://coopmaths.fr/alea/static/e3c/2020/tex/e3c_2020_00_sujet22_3_cor.tex
\begin{document}
\noindent\textit{APMEP / CoopMaths. CC BY-SA. Présentation adaptée par Maths Explorer.}\par\medskip
Dans un repère orthonormé, on considère le point $A(-3 ; 5)$ et la droite $(d)$ dont une équation cartésienne est $-x + 3y + 2 = 0$.
	
	\textbf{Question 1}
	
Avec l'équation de droite, on vérifie que les points $B(2; 0)$ et $C(-4; -2)$ appartenant à la droite $(d)$.

\begin{center}




	
	\begin{tikzpicture}[line cap=round,line join=round,>=stealth,x=.7cm,y=.7cm]
		
		\draw[step=1.0,gray, thin, dashed] (-5,-4) grid (6,3);
		
		
		\draw[->,color=black] (-5,0) -- (6,0);  
		\draw[->,color=black] (0,-4) -- (0,3);  
		
		
		\draw [fill=black] (0,0) circle (1.5pt);
		\draw[color=black] (0,0) node [anchor=north east]{$O$};
		
		
		\draw [->,line width=1.2pt,color=black] (0,0) -- (0,1);  
		\draw [->,line width=1.2pt,color=black] (0,0) -- (1,0);  
		
		
		\draw [color=black](0.5,0) node[anchor=north] {$\vec{\imath}$};  
		\draw [color=black](0,0.5) node[anchor=east] {$\vec{\jmath}$};   
		
		\draw [fill=black] (2,0) circle (1.5pt);
		\draw [fill=black] (2,0) node [anchor=north]{$B$};
			\draw [fill=black] (-4,-2) circle (1.5pt);
		\draw [fill=black] (-4,-2) node [anchor=south]{$C$};
		\draw[red,ultra thick](-5,-2.33)--(6,1.33);
			\draw [fill=black] (4,1) node [anchor=south]{$(d)$};
		
		\clip(-5,-5) rectangle (6,6);
	\end{tikzpicture}
	


\end{center}

	
	\textbf{Question 2}
	
	Un vecteur directeur de la droite $(d)$ a pour coordonnées $\begin{pmatrix}
		3\\
		1
	\end{pmatrix}$, donc un vecteur normal est par exemple $\vec{n}\begin{pmatrix}
	1\\
	-3
	\end{pmatrix}$. 
	
	\textbf{Question 3}
	
	Si $\delta$ est la perpendiculaire à $(d)$ passant par $A$, on a :
	\[
	M(x; y) \in \delta \iff \overrightarrow{AM} \text{ et } \vec{n} \text{ sont colinéaires, soit } -3(x + 3) = 1(y - 5).
	\]
	Cela donne l'équation :
	\[
	3x - y + 14 = 0.
	\]
	
\textbf{Question 4}

Le projeté orthogonal de $A$ sur la droite $(d)$ est le point commun de $(d)$ et $\delta$. On résout le système d'équations :
\[
\begin{cases}
	-x + 3y + 2 = 0 \\
	3x - y + 14 = 0
\end{cases}
\]
Cela donne :
\begin{align*}
	x &= 3y + 2 \\
	3(3y + 2) - y + 14 &= 0 \\
	9y + 6 - y + 14 &= 0 \\
	8y &= -20 \\
	y &= -\dfrac{5}{2}.
\end{align*}

On trouve $x = 3y + 2 = 3 \times \left(-\dfrac{5}{2}\right) + 2 = -\dfrac{11}{2}$.\\
 Le projeté orthogonal est donc le point $H\left(-\dfrac{11}{2}, -\dfrac{5}{2}\right)$.

\textbf{Question 5}

La distance entre le point $A$ et la droite $(d)$ est égale à $AH$. On calcule :
\begin{align*}
	AH^2 &= \left(-\dfrac{11}{2} + 3\right)^2 + \left(-\dfrac{5}{2} - 5\right)^2 \\
	&= \left(-\dfrac{5}{2}\right)^2 + \left(-\dfrac{15}{2}\right)^2 \\
	&= \dfrac{25}{4} + \dfrac{225}{4} \\
	&= \dfrac{250}{4}, \quad \text{d'où} \quad AH = \dfrac{\sqrt{250}}{2}.
\end{align*}

\end{document}
